% Source: https://artofproblemsolving.com/community/c3h1374100_preparation_contest
% Stable identity: aops-topic-1374100-q14
Use mass points. From $AF:FB=8:17$, masses $17$ at $A$ and $8$ at $B$ balance at $F$; from $AE:EC=2:5$, masses at $A$ and $C$ must be in the ratio $5:2$. Scaling, take
\[
m_A=85,\qquad m_B=40,\qquad m_C=34
\]
($85:40=17:8$ and $85:34=5:2$).

\begin{asy}
usepackage("amsmath");
size(5.5cm);
defaultpen(fontsize(10pt));
pair A=(1.5,4),B=(0,0),C=(6,0);
pair F=A+8/25*(B-A), Ep=A+2/7*(C-A), D=B+34/74*(C-B);
pair G=(85*A+40*B+34*C)/159;
draw(A--B--C--cycle,linewidth(0.9));
draw(A--D,blue); draw(B--Ep,blue); draw(C--F,blue);
dot(A);dot(B);dot(C);dot(D);dot(Ep);dot(F);dot(G,red);
label("$A$",A,N); label("$B$",B,SW); label("$C$",C,SE);
label("$D$",D,S); label("$E$",Ep,NE); label("$F$",F,W); label("$G$",G,2*dir(5),red);
label("$85$",A,2.4*E,deepgreen); label("$40$",B,3*NW,deepgreen); label("$34$",C,2.6*N,deepgreen);
label("\mbox{}",(min(currentpicture,user=true)+max(currentpicture,user=true))/2,black); // keeps the text after the figure black when LaTeX typesets the labels inline
\end{asy}

Then $G$ is the balance point of the whole system, $D$ carries mass $m_B+m_C$, and so on:
\[
\dfrac{AG}{GD}=\dfrac{m_B+m_C}{m_A}=\dfrac{74}{85},\qquad
\dfrac{BG}{GE}=\dfrac{m_A+m_C}{m_B}=\dfrac{119}{40},\qquad
\dfrac{CG}{GF}=\dfrac{m_A+m_B}{m_C}=\dfrac{125}{34}.
\]
The product is
\[
\dfrac{74\cdot119\cdot125}{85\cdot40\cdot34}=\dfrac{1100750}{115600}=\dfrac{1295}{136},
\]
which is in lowest terms ($1295=5\cdot7\cdot37$, $136=2^3\cdot17$). So $m+n=\boxed{1431}$.
